% SC1008 — Chapter 6: Inheritance, Polymorphism, Templates (0->100 ground-up)
\chapter{Inheritance, Polymorphism, and Templates}
\label{ch:inheritance-poly-templates}

\begin{quote}
\emph{``Inheritance is sharing; polymorphism is substituting; templates are stamping.''} Three ways to reuse code, each with a distinct cost model.
\end{quote}

\noindent\fbox{\parbox{\dimexpr\linewidth-2\fboxsep-2\fboxrule}{%
\textbf{From zero: reuse without copy-paste.} We start with why duplication hurts, then derive is-a hierarchies, virtual dispatch, and generic stamping from first principles.}}

\section*{Learning Outcomes}
After this chapter you will be able to:
\begin{enumerate}[label=(L\arabic*)]
    \item model is-a relationships with public inheritance;
    \item explain virtual dispatch and the vtable mechanism;
    \item use templates for generic functions and classes;
    \item distinguish static (compile-time) vs.\ dynamic (run-time) polymorphism;
    \item apply access specifiers and virtual destructors correctly.
\end{enumerate}

% ============================================================
\section{Inheritance: Is-A}
\label{sec:inheritance}
If \texttt{Dog} is-a \texttt{Animal}, every \texttt{Dog} has the \texttt{Animal} part plus its own extensions. Inheritance captures this.

\begin{figure}[h]\centering
\begin{tikzpicture}[scale=0.70,every node/.style={font=\scriptsize}]
\node[draw,thick,fill=blue!12,rounded corners=3pt,minimum width=3.0cm,minimum height=0.85cm,align=center] (animal) at (0,1.6) {\textbf{Animal}\\ \texttt{virtual void speak() const}\\ \texttt{virtual \~{}Animal()}};
\node[draw,thick,fill=green!14,rounded corners=3pt,minimum width=2.6cm,minimum height=0.85cm,align=center] (dog) at (-2.2,0) {\textbf{Dog : Animal}\\ \texttt{void speak() const override}\\ \texttt{void fetch()}};
\node[draw,thick,fill=orange!15,rounded corners=3pt,minimum width=2.6cm,minimum height=0.85cm,align=center] (cat) at (2.2,0) {\textbf{Cat : Animal}\\ \texttt{void speak() const override}};
\draw[-{Latex},thick,blue!60!black] (dog)--(animal);
\draw[-{Latex},thick,blue!60!black] (cat)--(animal);
\node[font=\tiny,fill=yellow!12,rounded corners=2pt,inner sep=2pt] at (0,-0.95) {arrow = is-a; derived inherits base subobject};
\end{tikzpicture}
\caption{Inheritance hierarchy: \texttt{Dog} and \texttt{Cat} each contain an \texttt{Animal} subobject. Public inheritance models is-a.}
\label{fig:inheritance}
\end{figure}

\begin{lstlisting}[language=C++]
struct Animal{
    virtual void speak() const { std::cout<<"...\n"; }
    virtual ~Animal() = default;           // virtual dtor: essential!
};
struct Dog : public Animal{
    void speak() const override { std::cout<<"Woof!\n"; }
    void fetch() const { std::cout<<"*fetches*\n"; }
};
Dog d; Animal& a = d; // Dog is-an Animal: binds base ref to derived
a.speak();            // Woof! (virtual dispatch)
\end{lstlisting}
\texttt{override} makes the compiler check the signature matches; always use it. Make base destructors \texttt{virtual} so \texttt{delete} through base pointer destroys the derived part.

\section{Polymorphism and Virtual Dispatch}
\label{sec:polymorphism}

\begin{figure}[h]\centering
\begin{tikzpicture}[scale=0.68,every node/.style={font=\scriptsize}]
% objects
\draw[thick,fill=blue!10] (0,0) rectangle (2.6,1.1);
\node[font=\tiny\bfseries] at (1.3,0.85) {\texttt{Dog d}};
\node[font=\tiny] at (1.3,0.55) {Animal part};
\node[font=\tiny] at (1.3,0.28) {Dog part};
\node[font=\tiny,fill=orange!18,rounded corners=1pt,inner sep=1pt] at (1.3,0.05) {vptr $\to$ Dog vtable};
\draw[thick,fill=green!10] (3.4,0) rectangle (6.0,1.1);
\node[font=\tiny\bfseries] at (4.7,0.85) {\texttt{Cat c}};
\node[font=\tiny] at (4.7,0.55) {Animal part};
\node[font=\tiny] at (4.7,0.28) {Cat part};
\node[font=\tiny,fill=orange!18,rounded corners=1pt,inner sep=1pt] at (4.7,0.05) {vptr $\to$ Cat vtable};
% vtables
\draw[thick,fill=violet!12] (0,-1.55) rectangle (2.6,-0.45);
\node[font=\tiny\bfseries] at (1.3,-0.60) {Dog vtable};
\node[font=\tiny] at (1.3,-0.92) {\texttt{speak $\to$ Dog::speak}};
\node[font=\tiny] at (1.3,-1.22) {\texttt{\~{}Animal $\to$ Dog::\~{}Dog}};
\draw[thick,fill=violet!12] (3.4,-1.55) rectangle (6.0,-0.45);
\node[font=\tiny\bfseries] at (4.7,-0.60) {Cat vtable};
\node[font=\tiny] at (4.7,-0.92) {\texttt{speak $\to$ Cat::speak}};
\draw[-{Latex},thick,red!60!black] (1.3,0.02)--(1.3,-0.45);
\draw[-{Latex},thick,red!60!black] (4.7,0.02)--(4.7,-0.45);
% call
\node[font=\tiny,fill=yellow!12,rounded corners=2pt,inner sep=2pt] at (3.0,-2.05) {\texttt{a.speak()} follows vptr $\to$ vtable $\to$ correct override};
\end{tikzpicture}
\caption{Virtual dispatch: each object holds a hidden \texttt{vptr} to its class's vtable. A virtual call indirects through it, selecting the derived override at run time.}
\label{fig:vtable}
\end{figure}

\begin{lstlisting}[language=C++]
void chorus(const std::vector<Animal*>& v){
    for(auto *a: v) a->speak(); // dynamic dispatch: Dog woofs, Cat meows
}
Dog d; Cat c;
chorus({&d,&c}); // Woof!  Meow!
\end{lstlisting}
Cost: one extra indirection per virtual call. Non-virtual calls are direct (zero overhead). Prefer non-virtual by default; make virtual only when you need run-time substitution. Pure virtual (\texttt{=0}) makes a base abstract: \texttt{virtual void draw() const =0;}.

\section{Templates: Compile-Time Stamping}
\label{sec:templates}
Templates generate code per type at compile time (static polymorphism). No vtable, no indirection --- but one instantiation per distinct type.

\begin{figure}[h]\centering
\begin{tikzpicture}[scale=0.70,every node/.style={font=\scriptsize}]
\node[draw,thick,fill=blue!10,rounded corners=3pt,minimum width=3.4cm,minimum height=0.85cm,align=center] (tmpl) at (0,0) {\texttt{template<typename T>}\\ \texttt{T max2(T a,T b)}};
\node[draw,thick,fill=green!14,rounded corners=3pt,minimum width=2.4cm,minimum height=0.70cm] (i) at (4.8,0.55) {\texttt{max2<int>}};
\node[draw,thick,fill=orange!15,rounded corners=3pt,minimum width=2.4cm,minimum height=0.70cm] (d) at (4.8,-0.55) {\texttt{max2<double>}};
\draw[-{Latex},thick,blue!60!black] (tmpl)--(i) node[midway,above,font=\tiny] {stamp};
\draw[-{Latex},thick,blue!60!black] (tmpl)--(d) node[midway,below,font=\tiny] {stamp};
\node[font=\tiny,fill=yellow!12,rounded corners=2pt,inner sep=2pt] at (0,-1.10) {one template, many generated functions};
\end{tikzpicture}
\caption{Templates as stamps: the compiler generates a separate function/class per \texttt{T}. No run-time cost, but code bloat if overused.}
\label{fig:template-stamp}
\end{figure}

\begin{lstlisting}[language=C++]
template<typename T>
T max2(T a, T b){ return (a>b)?a:b; }
int x = max2<int>(3,7);          // explicit
double y = max2(3.1, 2.9);       // deduced: T=double
// Class template
template<typename T>
struct VecT{ T x,y; T dot(const VecT& o) const { return x*o.x + y*o.y; } };
VecT<double> v{1.0,2.0};
\end{lstlisting}
Templates are type-checked after instantiation; errors can be verbose. C++20 \texttt{concept}s constrain \texttt{T} with readable diagnostics (preview in Ch.~8).

\section{Static vs.\ Dynamic Polymorphism}
\begin{center}\small
\begin{tabular}{lcc}
\toprule & Templates (static) & Virtual (dynamic)\\
\midrule When chosen & compile time & run time\\
Cost & code size per \texttt{T} & one vtable indirection\\
Heterogeneous container & no (one \texttt{T} per container) & yes (\texttt{vector<Animal*>})\\
\bottomrule
\end{tabular}
\end{center}
Rule of thumb: use templates for generic algorithms on uniform types; use virtual for heterogeneous runtime substitution.


\section{Access Control and Protected Members}
\texttt{public} inheritance keeps is-a; \texttt{private}/\texttt{protected} inheritance is implementation reuse, not is-a. \texttt{protected} members are visible to derived classes but not to outsiders.

\begin{lstlisting}[language=C++]
class Base{
protected: int prot_ = 42;  // derived can see, outsiders cannot
private:   int priv_ = 0;
public:    int pub_  = 1;
};
class Derived : public Base{
public:
    void foo(){ prot_ = 7; /* ok */ /* priv_=1; error */ }
};
Base b; // b.prot_ error -- protected not accessible outside
\end{lstlisting}
Prefer \texttt{private} data with \texttt{protected} accessors over \texttt{protected} data: it preserves the base's invariant while still allowing customisation.

\section{Abstract Classes and Interfaces}
A class with at least one pure virtual function (\texttt{=0}) is abstract --- it cannot be instantiated, only derived from. This models interfaces.

\begin{lstlisting}[language=C++]
struct Drawable{ // interface: pure virtual
    virtual void draw() const = 0;
    virtual ~Drawable() = default;
};
struct Circle : Drawable{
    double r; explicit Circle(double r_): r(r_) {}
    void draw() const override { std::cout<<"Circle r="<<r<<"\n"; }
};
struct Square : Drawable{
    double s; explicit Square(double s_): s(s_) {}
    void draw() const override { std::cout<<"Square s="<<s<<"\n"; }
};
void render(const std::vector<Drawable*>& v){
    for(auto *d: v) d->draw(); // uniform interface
}
\end{lstlisting}
Interfaces decouple callers from concrete types --- the foundation of plugin architectures and test mocks. In modern C++, pure abstract classes are preferred over deep hierarchies.


\section{Template Specialisation and Non-Type Parameters}
Templates can be specialised for specific types or take compile-time constants as parameters.

\begin{lstlisting}[language=C++]
template<typename T, int N> // non-type param: array size known at compile time
struct Array{
    T data[N];
    T& operator[](int i){ return data[i]; }
    constexpr int size() const { return N; }
};
Array<int,5> a; // 5 ints, no heap, size() is compile-time

// Function template with overload for const char*
template<typename T> void print2(T x){ std::cout<<x<<"\n"; }
template<> void print2<const char*>(const char* s){ std::cout<<"C-string: "<<s<<"\n"; }
print2(42);        // int version
print2("hello");   // specialised for const char*
\end{lstlisting}
Non-type parameters enable fixed-size containers without dynamic allocation; specialisation lets you tailor behaviour per type while keeping a uniform call site.

\section{Chapter Notes}
Sutter \& Alexandrescu \emph{C++ Coding Standards} Ch.~37--40 covers inheritance discipline. STL (Ch.~7) is almost entirely templates; inheritance appears sparingly there by design.

\section{Exercises}
\begin{enumerate}[label=E6.\arabic*, leftmargin=*, itemsep=0.6em]
\item \textbf{Slicing.} \texttt{Animal a = Dog(); a.speak();} What prints and why? How does \texttt{Animal\& a = d;} differ? Explain object slicing.
\item \textbf{Virtual dtor.} Give a program where deleting a \texttt{Dog} through \texttt{Animal*} without a virtual destructor leaks resources. Fix it.
\item \textbf{Template deduction.} For \texttt{template<typename T> void f(T a, T b)}, what happens with \texttt{f(1, 2.0)}? How to fix the call or the template?
\item \textbf{Override check.} Remove \texttt{override} and misspell \texttt{speak} as \texttt{speal} in \texttt{Dog}. What happens with and without \texttt{override}? Why is \texttt{override} best practice?
\item \textbf{Design.} Model \texttt{Shape} with \texttt{area()} and \texttt{draw()}. Should they be virtual/pure virtual? Should \texttt{Shape} be instantiable? Justify and sketch the hierarchy.
\end{enumerate}
% End of Chapter 6
